% =============================================================================
% bab1-pendahuluan-pb.tex — BAB I PENDAHULUAN (Proyek Bisnis / Studi Kasus)
% Lampiran II — S2 only
% 6 Sub-bab (berbeda dengan Lampiran I yang 5 sub-bab)
% =============================================================================

\chapter{PENDAHULUAN}
\label{bab:pendahuluan-pb}

\section{Gambaran Umum Objek Penelitian}
\label{sec:gambaran-umum-pb}

Uraikan gambaran umum perusahaan/organisasi/industri yang menjadi objek studi.
Jelaskan profil singkat, sejarah, visi-misi, dan kondisi bisnis terkini dari
objek penelitian.

\section{Latar Belakang Penelitian}
\label{sec:latar-belakang-pb}

Uraikan latar belakang permasalahan bisnis yang dihadapi objek penelitian.
Jelaskan fenomena empiris, kondisi internal/eksternal organisasi, serta justifikasi
pentingnya penelitian ini dilakukan untuk menyelesaikan masalah tersebut.

\section{Rumusan Masalah}
\label{sec:rumusan-masalah-pb}

Berdasarkan uraian latar belakang di atas, rumusan masalah dalam penelitian ini adalah:

\begin{enumerate}
  \item Bagaimana kondisi [aspek bisnis tertentu] pada [nama organisasi]?
  \item Apa saja faktor yang mempengaruhi [permasalahan bisnis]?
  \item Bagaimana solusi/rekomendasi yang tepat untuk mengatasi [permasalahan bisnis]?
\end{enumerate}

\section{Tujuan Penelitian}
\label{sec:tujuan-pb}

Tujuan dari penelitian ini adalah:

\begin{enumerate}
  \item Untuk menganalisis kondisi [aspek bisnis tertentu] pada [nama organisasi].
  \item Untuk mengidentifikasi faktor-faktor yang mempengaruhi [permasalahan bisnis].
  \item Untuk merumuskan rekomendasi solusi atas [permasalahan bisnis] pada [nama organisasi].
\end{enumerate}

\section{Manfaat Penelitian}
\label{sec:manfaat-pb}

Hasil penelitian ini diharapkan dapat memberikan manfaat bagi:

\begin{enumerate}
  \item \textbf{Manfaat Teoritis:} Memperkaya khazanah ilmu pengetahuan di bidang
        [bidang ilmu] khususnya mengenai [topik penelitian].

  \item \textbf{Manfaat Praktis:} Memberikan rekomendasi solusi yang dapat diterapkan
        oleh [nama organisasi] dalam mengatasi [permasalahan bisnis].
\end{enumerate}

\section{Sistematika Penulisan}
\label{sec:sistematika-pb}

Penulisan proyek bisnis ini disusun dalam lima bab dengan sistematika sebagai berikut:

\textbf{BAB I PENDAHULUAN}, berisi gambaran umum objek penelitian, latar belakang
penelitian, rumusan masalah, tujuan penelitian, manfaat penelitian, dan sistematika
penulisan.

\textbf{BAB II TINJAUAN PUSTAKA ASPEK BISNIS}, berisi tinjauan teori yang relevan,
kerangka pemikiran, dan kajian penelitian terdahulu.

\textbf{BAB III METODE ANALISIS BISNIS}, berisi metode penelitian, jenis dan sumber
data, teknik pengumpulan data, serta teknik analisis data.

\textbf{BAB IV IMPLEMENTASI PROYEK BISNIS/STUDI KASUS}, berisi hasil analisis dan
implementasi solusi atas permasalahan bisnis yang diteliti.

\textbf{BAB V KESIMPULAN DAN SARAN}, berisi kesimpulan dari hasil penelitian dan
saran/rekomendasi bagi pihak-pihak yang berkepentingan.
